\editionPage[\pdfbookmark[0]{Indications des défis}{bm-indices-debut-parent}]{indices-debut}{Indications progressives}{Défis 1 à 3}{Ne lire que l’aide nécessaire.}
\begin{blocEdition}{indices-debut}
\editionEtape{D01 · Les nombres miroirs}\textbf{Indice 1.} Le premier chiffre est 4 : que vaut alors le dernier ?\par\medskip\textbf{Indice 2.} Seuls le deuxième et le troisième chiffre sont libres.\par
\editionEtape{D02 · Somme et écart}\textbf{Indice 1.} Si on retire l’écart à la somme, que reste-t-il ?\par\medskip\textbf{Indice 2.} Ce qui reste se partage en deux parts égales.\par
\editionEtape{D03 · Les élèves en rangs}\textbf{Indice 1.} Que se passe-t-il si un élève de plus arrive ?\par\medskip\textbf{Indice 2.} Avec un élève de plus, le nombre est multiple de 4, de 5 et de 6.\par

\end{blocEdition}
\editionPage{indices-4}{Indications progressives}{Défis 4 à 6}{Ne lire que l’aide nécessaire.}
\begin{blocEdition}{indices-4}
\editionEtape{D04 · Les divisions qui s’arrêtent}\textbf{Indice 1.} Poser $1:8$ : 10 dixièmes, puis les centièmes, puis les millièmes.\par\medskip\textbf{Indice 2.} Pour 3, 6, 7 et 9, observer les restes successifs.\par
\editionEtape{D05 · La suite des lapins}\textbf{Indice 1.} Écrire les termes un à un en additionnant les deux derniers.\par\medskip\textbf{Indice 2.} Continuer jusqu’à dépasser 1 000.\par
\editionEtape{D06 · La fourmi pressée}\textbf{Indice 1.} Après un saut, elle est à la moitié du mètre ; il reste une moitié.\par\medskip\textbf{Indice 2.} À chaque saut, ce qui reste est divisé par deux.\par

\end{blocEdition}
\editionPage{indices-7}{Indications progressives}{Défis 7 à 9}{Ne lire que l’aide nécessaire.}
\begin{blocEdition}{indices-7}
\editionEtape{D07 · Les douzièmes du milieu}\textbf{Indice 1.} Écrire $\frac13$ et $\frac34$ en douzièmes.\par\medskip\textbf{Indice 2.} Lister les numérateurs strictement entre les deux bornes.\par
\editionEtape{D08 · L’héritage des chameaux}\textbf{Indice 1.} Calculer la moitié, le tiers et le neuvième de 18.\par\medskip\textbf{Indice 2.} Additionner $\frac12+\frac13+\frac19$ en dix-huitièmes.\par
\editionEtape{D09 · Le décimal à code}\textbf{Indice 1.} Appeler d le chiffre des dixièmes : que vaut celui des millièmes ?\par\medskip\textbf{Indice 2.} La somme des chiffres décimaux fait trois fois d.\par

\end{blocEdition}
\editionPage{indices-10}{Indications progressives}{Défis 10 à 12}{Ne lire que l’aide nécessaire.}
\begin{blocEdition}{indices-10}
\editionEtape{D10 · La tirelire à deux pièces}\textbf{Indice 1.} Si les 10 pièces étaient toutes de 20 centimes, quel serait le total ?\par\medskip\textbf{Indice 2.} Remplacer une pièce de 20 c par une de 50 c ajoute 30 centimes.\par
\editionEtape{D11 · La balance aux trois masses}\textbf{Indice 1.} Une masse posée à côté de l’objet se retranche.\par\medskip\textbf{Indice 2.} Pour 2 g : comparer l’objet plus 1 g à la masse de 3 g.\par
\editionEtape{D12 · Le plus grand enclos}\textbf{Indice 1.} Longueur plus largeur vaut la moitié du périmètre.\par\medskip\textbf{Indice 2.} Calculer l’aire de chaque paire possible.\par

\end{blocEdition}
\editionPage{indices-13}{Indications progressives}{Défis 13 à 15}{Ne lire que l’aide nécessaire.}
\begin{blocEdition}{indices-13}
\editionEtape{D13 · L’horloge des angles}\textbf{Indice 1.} L’angle droit à 3 h correspond à trois « heures » du cadran.\par\medskip\textbf{Indice 2.} Une heure du cadran correspond donc à $90:3$ degrés.\par
\editionEtape{D14 · Les deux sabliers}\textbf{Indice 1.} Retourner les deux sabliers ensemble au départ.\par\medskip\textbf{Indice 2.} À 8 minutes, il reste exactement 1 minute de sable dans le bas du sablier de 7 minutes.\par
\editionEtape{D15 · Le trésor à deux distances}\textbf{Indice 1.} Où sont tous les points à 4 cm du puits ?\par\medskip\textbf{Indice 2.} Chercher les points communs à deux cercles.\par

\end{blocEdition}
\editionPage{indices-16}{Indications progressives}{Défis 16 à 18}{Ne lire que l’aide nécessaire.}
\begin{blocEdition}{indices-16}
\editionEtape{D16 · Les diagonales des polygones}\textbf{Indice 1.} Depuis un sommet, on ne peut tracer de diagonale ni vers lui-même ni vers ses deux voisins.\par\medskip\textbf{Indice 2.} Chaque diagonale est comptée deux fois, une fois par extrémité.\par
\editionEtape{D17 · Les mots au fil de l’eau}\textbf{Indice 1.} Un reflet horizontal retourne chaque lettre de haut en bas, sans changer leur ordre.\par\medskip\textbf{Indice 2.} Il suffit que toutes les lettres aient un axe horizontal.\par
\editionEtape{D18 · Le grand cube peint}\textbf{Indice 1.} Les cubes à 3 faces peintes sont aux coins.\par\medskip\textbf{Indice 2.} Les cubes sans peinture forment un cube intérieur de 2 sur 2 sur 2.\par

\end{blocEdition}
\editionPage{indices-19}{Indications progressives}{Défis 19 à 21}{Ne lire que l’aide nécessaire.}
\begin{blocEdition}{indices-19}
\editionEtape{D19 · Deux robinets ensemble}\textbf{Indice 1.} Combien de litres chaque robinet donne-t-il en une minute ?\par\medskip\textbf{Indice 2.} Ensemble, ils additionnent leurs débits.\par
\editionEtape{D20 · La récréation en secteurs}\textbf{Indice 1.} Calculer d’abord le nombre de joueurs de ballon et de corde.\par\medskip\textbf{Indice 2.} Le reste se compare au total de 120.\par
\editionEtape{D21 · La somme la plus probable}\textbf{Indice 1.} Construire le tableau à double entrée des 36 issues.\par\medskip\textbf{Indice 2.} Compter les cases de chaque somme.\par

\end{blocEdition}
